% ============================================================
\chapter{Modul 1: Einstieg in Python}
% ============================================================

Python ist heute eine der meistgenutzten Programmiersprachen in der Robotik -- von einfachen Microcontroller-Skripten bis zu komplexen KI-Systemen. Dieses Modul legt das Fundament.

% ────────────────────────────────────────────────────────────
\section{Lektion 1: Was ist Python?}
% ────────────────────────────────────────────────────────────

\subsection*{Lernziele}
\begin{itemize}
    \item verstehen, was eine Programmiersprache ist,
    \item erklären, warum Python für Robotik geeignet ist,
    \item das erste Python-Programm schreiben und verstehen.
\end{itemize}

\subsection*{Wie denkt ein Computer?}

\begin{analogie}
Stell dir einen sehr präzisen, aber völlig ahnungslosen Assistenten vor. Du kannst ihm alles beibringen -- aber nur in exakten, eindeutigen Schritten. Sagst du ``Mach das Licht an'', reagiert er nicht. Sagst du ``Gehe zum Schalter, drücke ihn nach oben'', funktioniert es.

Genau so arbeitet ein Computer. Er ist blitzschnell und unermüdlich -- aber er versteht nur präzise Anweisungen.
\end{analogie}

Eine \textbf{Programmiersprache} ist das Werkzeug, mit dem wir solche Anweisungen formulieren. Python ist dabei besonders gut lesbar -- Python-Code sieht aus wie vereinfachtes Englisch.

\subsection*{Warum Python für Robotik?}

\begin{praxis}
ROS2 (Robot Operating System 2), das wichtigste Framework für professionelle Robotik, unterstützt Python als erste Sprache. OpenCV (Bildverarbeitung), NumPy (Mathematik), scikit-learn (Machine Learning) -- alle haben Python als Hauptsprache. Wer in der Robotik arbeiten will, kommt an Python nicht vorbei.
\end{praxis}

\subsection*{Das erste Programm}

\begin{lstlisting}
print("Roboter bereit.")
print("Warte auf Befehl...")
\end{lstlisting}

\begin{merksatz}
\texttt{print()} gibt Text auf dem Bildschirm aus. Der Text steht in Anführungszeichen -- das teilt Python mit: Das ist ein Text, keine Berechnung.
\end{merksatz}

\begin{uebung}[title=Übung 1.1: Erste Ausgaben]
Schreibe ein Programm, das drei Zeilen ausgibt:
\begin{itemize}
    \item den Namen deines Roboters,
    \item seine Funktion (z.\,B. ``Patrouille''),
    \item den Status: ``Bereit.''
\end{itemize}
\end{uebung}

\begin{ros}
In ROS2 gibt jeder Node beim Start eine Meldung aus. Das passiert intern mit \texttt{self.get\_logger().info("Node gestartet")} -- im Prinzip ein erweitertes \texttt{print()}.
\end{ros}

% ────────────────────────────────────────────────────────────
\section{Lektion 2: Variablen}
% ────────────────────────────────────────────────────────────

\subsection*{Lernziele}
\begin{itemize}
    \item erklären, was eine Variable ist,
    \item Variablen anlegen und ausgeben,
    \item sinnvolle Variablennamen vergeben.
\end{itemize}

\begin{analogie}
Eine Variable ist wie eine beschriftete Schublade. Auf dem Zettel steht der Name (\texttt{geschwindigkeit}), in der Schublade liegt der Wert (\texttt{0.3}). Der Name bleibt -- der Inhalt kann sich ändern.
\end{analogie}

\begin{lstlisting}
roboter_name   = "MobileRobot-1"
geschwindigkeit = 0.3          # Meter pro Sekunde
akku_prozent    = 87
motor_aktiv     = True

print(roboter_name)
print(f"Geschwindigkeit: {geschwindigkeit} m/s")
print(f"Akku: {akku_prozent}%")
\end{lstlisting}

\begin{merksatz}
\textbf{Benennungsregeln:}
\begin{itemize}
    \item Nur Buchstaben, Zahlen, Unterstriche (\texttt{\_})
    \item Keine Umlaute, kein Leerzeichen
    \item Aussagekräftig: \texttt{akku\_prozent} statt \texttt{a}
    \item Für mehrere Wörter: \texttt{snake\_case} (Python-Standard)
\end{itemize}
\end{merksatz}

\subsection*{f-Strings: Variablen in Text einbetten}

\begin{lstlisting}
name = "MobileRobot-1"
akku = 87

# Klassisch (umständlich):
print("Roboter " + name + " hat " + str(akku) + "% Akku.")

# Mit f-String (modern, empfohlen):
print(f"Roboter {name} hat {akku}% Akku.")
\end{lstlisting}

\begin{praxis}
Statusmeldungen in Robotersystemen verwenden intensiv f-Strings: Position, Sensor-wert, Timestamp -- alles wird in einem lesbaren String zusammengesetzt.
\end{praxis}

\begin{uebung}[title=Übung 1.2: Roboter-Datensatz]
Lege Variablen für einen Roboter an: Name, maximale Geschwindigkeit, Akkustand, Anzahl Räder. Gib alle Werte in einem übersichtlichen f-String aus.
\end{uebung}

% ────────────────────────────────────────────────────────────
\section{Lektion 3: Datentypen}
% ────────────────────────────────────────────────────────────

\subsection*{Lernziele}
\begin{itemize}
    \item die vier Grunddatentypen kennen und unterscheiden,
    \item \texttt{type()} zur Typprüfung nutzen,
    \item Typen konvertieren.
\end{itemize}

Python unterscheidet automatisch, welche Art von Daten eine Variable enthält.

\begin{center}
\begin{tabular}{llll}
\toprule
\textbf{Typ} & \textbf{Beispiel} & \textbf{Robotik-Anwendung} & \textbf{Python} \\
\midrule
Ganzzahl   & 5, -3, 200     & Schritte, Impulse, Zähler  & \texttt{int}   \\
Kommazahl  & 0.35, -9.81    & Sensorwerte, Koordinaten   & \texttt{float} \\
Text       & "links", "OK"  & Befehle, Statusmeldungen   & \texttt{str}   \\
Wahrheit   & True, False    & Motor an/aus, Sensor OK    & \texttt{bool}  \\
\bottomrule
\end{tabular}
\end{center}

\begin{lstlisting}
# Sensorpaket eines mobilen Roboters
entfernung  = 0.42       # float: Meter
winkel      = 90         # int: Grad
richtung    = "links"    # str
notaus      = False      # bool

# Typ prüfen
print(type(entfernung))   # <class 'float'>
print(type(winkel))       # <class 'int'>
\end{lstlisting}

\begin{merksatz}
\texttt{sensor\_wert = "0.42"} und \texttt{sensor\_wert = 0.42} sind NICHT dasselbe. Ersteres ist Text -- damit kann man nicht rechnen. Das ist ein sehr häufiger Fehler bei Sensor-Eingaben.
\end{merksatz}

\subsection*{Typkonvertierung}

\begin{lstlisting}
# Benutzereingabe ist immer ein String
eingabe = "150"         # kommt von input() oder Serial
schritte = int(eingabe) # jetzt kann man rechnen

# Zahl in Text
akku = 87
nachricht = "Akku: " + str(akku) + "%"

# int zu float und umgekehrt
x = float(3)    # 3.0
y = int(3.9)    # 3  (abschneiden, nicht runden!)
\end{lstlisting}

\begin{praxis}
Serielle Kommunikation (Arduino zu Python über USB) überträgt immer Strings. Das erste, was ein Robotik-Programm tut: Strings in Zahlen konvertieren, damit es rechnen kann.
\end{praxis}

\begin{uebung}[title=Übung 1.3: Typen bestimmen und konvertieren]
Gegeben:
\begin{lstlisting}
a = "3.14"
b = "42"
c = 7
\end{lstlisting}
Konvertiere \texttt{a} zu \texttt{float}, \texttt{b} zu \texttt{int}, \texttt{c} zu \texttt{str} und gib alle Typen mit \texttt{type()} aus.
\end{uebung}

% ────────────────────────────────────────────────────────────
\section{Lektion 4: Rechnen mit Python}
% ────────────────────────────────────────────────────────────

\subsection*{Lernziele}
\begin{itemize}
    \item alle Rechenoperatoren kennen,
    \item Formeln aus der Robotik in Python umsetzen,
    \item Operatorrangfolge verstehen.
\end{itemize}

\begin{lstlisting}
# Kinematik: Weg-Zeit-Rechnung
v = 0.3     # Geschwindigkeit in m/s
t = 5.0     # Zeit in Sekunden
s = v * t   # Weg in Metern
print(f"Zurueckgelegter Weg: {s} m")   # 1.5 m

# Ohmsches Gesetz: Strom durch Motorwicklung
U = 12.0    # Volt
R = 240     # Ohm
I = U / R   # Ampere
print(f"Strom: {I:.4f} A")   # 0.0500 A

# Modulo: Encoder-Tick-Ueberlauf
ticks      = 1250
ticks_max  = 1024
normiert   = ticks % ticks_max
print(f"Normierter Tickwert: {normiert}")   # 226

# Potenz: Pythagoras fuer Abstandsberechnung
dx, dy = 3.0, 4.0
d = (dx**2 + dy**2) ** 0.5
print(f"Euklidischer Abstand: {d} m")   # 5.0
\end{lstlisting}

\begin{merksatz}
Die sieben Operatoren: \texttt{+} \texttt{-} \texttt{*} \texttt{/} (float-Division) \texttt{//} (Ganzzahl-Division) \texttt{\%} (Rest) \texttt{**} (Potenz). Rangfolge wie in der Mathematik -- Klammern im Zweifel immer setzen.
\end{merksatz}

\begin{praxis}
Kinematik, Sensorverarbeitung, Regelungstechnik -- alles in der Robotik ist Mathematik. Python ersetzt den Taschenrechner und macht Berechnungen wiederholbar und automatisierbar.
\end{praxis}

\begin{uebung}[title=Projektaufgabe Modul 1: Roboter-Statusbericht]
Schreibe ein vollständiges Python-Programm für deinen Roboter. Definiere:
\begin{itemize}
    \item Name, Geschwindigkeit (\texttt{float}), Akkustand (\texttt{int}), Motor-Status (\texttt{bool})
    \item Berechne: Reichweite = Akku-Energie / Leistung (nimm sinnvolle Werte)
    \item Gib einen formatierten Statusbericht aus (mindestens 5 Zeilen)
\end{itemize}
\end{uebung}
